% Section 4: Channel 1 -- Coordination Failure
\label{sec:channel1}

\subsection{Setup}

We embed the $\rho$-parameterised signal structure of Section~\ref{sec:model-signal} into the \citet{goldstein2005} bank-run framework. A continuum of depositors indexed by $i \in [0,1]$ each choose whether to withdraw ($a_i = 1$) or keep funds ($a_i = 0$). The bank holds assets with quality $\theta \sim \mathrm{Uniform}[0,1]$ and fails if the aggregate measure of withdrawers $\ell(a) \equiv \int a_i \, di$ exceeds the fundamental: $\ell(a) > \theta$. Payoffs follow the standard \citet{goldstein2005} structure. A withdrawing agent receives one unit if the bank survives and a pro-rata share $\theta/\ell(a)$ if it fails. An agent who keeps funds receives $R(\theta) > 1$ if the bank survives and zero if it fails, where $R$ is continuous and strictly increasing with $R(0) < 1 < R(1)$.

A fraction $\lambda$ of depositors are AI-equipped and observe signals $x_i = \theta + \sqrt{\rho}\,\eta + \sqrt{1-\rho}\,\xi_i$ as in~\eqref{eq:noise-decomposition}. The remaining fraction $1-\lambda$ are traditional agents who observe $x_i = \theta + \xi_i$. We seek a symmetric monotone Bayesian Nash equilibrium in threshold strategies: each agent withdraws if and only if $x_i < x^*$.

The key departure from standard \citet{goldstein2005} is that the aggregate withdrawal fraction $\ell$ depends on the common shock $\eta$. For AI-equipped agents, the fraction withdrawing given threshold $x^*$, fundamental $\theta$, and shock $\eta$ is
\begin{equation}\label{eq:withdrawal-ai}
  \ell_{\text{AI}}(\theta, \eta, x^*) = \Phi\!\left(\frac{x^* - \theta - \sqrt{\rho}\,\eta}{\sqrt{1-\rho}\,\sigma}\right),
\end{equation}
where $\Phi$ denotes the standard normal CDF. The AI withdrawal fraction in~\eqref{eq:withdrawal-ai} depends on $\eta$ through $\sqrt{\rho}$, creating correlated withdrawal. For traditional agents it is $\ell_T(\theta, x^*) = \Phi((x^*-\theta)/\sigma)$. The total withdrawal fraction is
\begin{equation}\label{eq:withdrawal-total}
  \ell(\theta, \eta, x^*) = \lambda\,\ell_{\text{AI}}(\theta, \eta, x^*) + (1-\lambda)\,\ell_T(\theta, x^*).
\end{equation}
The dependence of $\ell$ on $\eta$ means that AI-equipped agents withdraw in a correlated pattern. The crisis threshold $\theta^*(\rho)$ is defined by the condition that the expected aggregate withdrawal equals the fundamental:
\begin{equation}\label{eq:crisis-threshold}
  \mathbb{E}_\eta\!\left[\ell(\theta^*, \eta, x^*)\right] = \theta^*.
\end{equation}

\subsection{Equilibrium Characterisation}

Each AI-equipped agent's signal $x_i$ can be decomposed as
\begin{equation}\label{eq:effective-fundamental}
  x_i = \underbrace{\theta + \sqrt{\rho}\,\eta}_{\tilde{\theta}} + \sqrt{1-\rho}\,\xi_i,
\end{equation}
where $\tilde{\theta} \equiv \theta + \sqrt{\rho}\,\eta$ is the effective fundamental. The decomposition in~\eqref{eq:effective-fundamental} shows that the common shock $\sqrt{\rho}\,\eta$ shifts all AI-equipped agents' signals identically, functioning as a public signal in the sense of \citet{hellwig2002}: it creates correlated posterior beliefs and, through them, correlated withdrawal decisions.

The Morris-Shin uniqueness condition requires that effective public precision remain below a threshold relative to private precision. Here, the effective public precision is $\tau_{\text{pub}} = 1/(\rho\sigma^2)$ and the private precision is $\tau_{\text{priv}} = 1/((1-\rho)\sigma^2)$. Applying the \citet{morris2003} condition $\alpha_{\text{SC}} \cdot \sqrt{\rho/(1-\rho)} < 1$, where $\alpha_{\text{SC}} \in (0,1)$ is the strategic complementarity parameter, yields the critical correlation

\begin{equation}\label{eq:rho1star}
  \rho_1^* = \frac{1}{1 + \alpha_{\text{SC}}^2}.
\end{equation}
For $\rho < \rho_1^*$ in~\eqref{eq:rho1star}, the uniqueness condition holds and a unique threshold equilibrium exists. For $\rho \geq \rho_1^*$, the \citet{hellwig2002} multiplicity region is restored: both a bank-run equilibrium and a no-run equilibrium coexist for fundamentals in a neighbourhood of $\theta^*$.

\begin{proposition}[Uniqueness-Multiplicity Boundary]\label{prop:uniqueness-boundary}
  In the Goldstein-Pauzner (2005) bank-run game with the $\rho$-parameterised signal structure~\eqref{eq:noise-decomposition} and a fraction $\lambda$ of AI-equipped agents, the threshold equilibrium is unique if and only if $\rho < \rho_1^*$, where $\rho_1^*$ is defined in~\eqref{eq:rho1star}.
  For $\rho > \rho_1^*$, the global game admits multiple equilibria and the Morris-Shin uniqueness selection breaks down.

  \prooflabel{Proof sketch.} The Goldstein-Pauzner dominance regions ensure monotone threshold strategies (\citealt{goldstein2005}, Lemma~1), reducing the equilibrium to a fixed-point system $(F_1, F_2)$ in $(x^*, \theta^*)$. The Jacobian determinant of this system depends on the common-to-private precision ratio $\alpha_{\text{SC}} \cdot \sqrt{\rho/(1-\rho)}$, exactly as in the linear-quadratic case of \citet{morris2003}.

  \prooflabel{Assumption~A1.} We assume the uniqueness boundary for the binary-action game is governed by the same critical ratio: uniqueness holds if and only if $\alpha_{\text{SC}} \cdot \sqrt{\rho/(1-\rho)} < 1$. This is consistent with \citet{hellwig2002} (Theorem~1), which requires only monotone best responses and the dominance-region property, both satisfied here. Equivalently, $\rho < 1/(1+\alpha_{\text{SC}}^2)$. $\square$
\end{proposition}

The formula~\eqref{eq:rho1star} has the correct limiting behaviour. When $\alpha_{\text{SC}} = 0$ (no strategic complementarity), $\rho_1^* = 1$: uniqueness always holds regardless of signal correlation. When $\alpha_{\text{SC}} \to \infty$ (extreme complementarity), $\rho_1^* \to 0$: even infinitesimal correlation destroys uniqueness. For the intermediate case $\alpha_{\text{SC}} = 1$, $\rho_1^* = 0.5$.

\subsection{Non-Monotonicity of the Crisis Threshold}

\begin{proposition}[Crisis Threshold Non-Monotonicity]\label{prop:crisis-nonmonotone}
  Consider the Goldstein-Pauzner (2005) bank-run game with $\rho$-parameterised AI signals and noise parameter $\sigma$ satisfying $\sigma < (\theta_H - \theta_L)\sqrt{1-\rho_1^*}/(2\sqrt{2\pi})$, which ensures the global games posterior concentrates in the interior of the dominance region for all $\rho \in [0, \rho_1^*)$. For $\rho \in (0, \rho_1^*)$, the unique crisis threshold $\theta^*(\rho)$ is non-monotone in $\rho$: there exists $\rho_0 \in (0, \rho_1^*)$ such that
  \begin{equation}\label{eq:theta-star-monotonicity}
    \frac{d\theta^*}{d\rho} < 0 \quad \text{for } \rho \in (0, \rho_0), \qquad
    \frac{d\theta^*}{d\rho} > 0 \quad \text{for } \rho \in (\rho_0, \rho_1^*).
  \end{equation}
  The comparative statics of $\theta^*(\rho)$ are qualitatively distinct from those of the Danielsson-Uthemann (2025) threshold $\theta^*(\mu) = (c/b)[(1-\mu)p + \mu]$, which is monotonically increasing and linear in $\mu$.

  \prooflabel{Proof sketch.} Two competing effects determine the sign of $d\theta^*/d\rho$. (a) \emph{Precision effect:} as $\rho$ increases, the idiosyncratic noise variance $(1-\rho)\sigma^2$ falls, making the private signal component more precise. Better individual inference reduces coordination failure, lowering $\theta^*$. (b) \emph{Common noise effect:} as $\rho$ increases, the common variance $\rho\sigma^2$ rises. The common shock $\eta$ shifts all AI-equipped agents' signals identically, increasing the correlation of withdrawal decisions conditional on $\theta$ and making coordinated bank runs easier to sustain, raising $\theta^*$. At low $\rho$, effect (a) dominates; at high $\rho$, effect (b) dominates. The turning point $\rho_0$ is found by differentiating the implicit system~\eqref{eq:crisis-threshold} and~\eqref{eq:withdrawal-total} and setting $d\theta^*/d\rho = 0$. By the implicit function theorem applied to the Goldstein-Pauzner indifference condition, the derivative exists and is continuous for $\rho < \rho_1^*$. $\square$
\end{proposition}

The regularity condition on $\sigma$ in Proposition~\ref{prop:crisis-nonmonotone} is sufficient but not necessary. It ensures that the posterior beliefs concentrate in the interior of the dominance region, so that the implicit function theorem applies to the Goldstein-Pauzner system for all $\rho \in [0, \rho_1^*)$. The non-monotonicity result may hold under weaker conditions, but the sufficient condition guarantees the existence and continuity of $d\theta^*/d\rho$ needed for the proof.

The sign reversal in~\eqref{eq:theta-star-monotonicity} has a sharp economic interpretation. At low $\rho$, correlated AI signals function like better private information: agents are individually more accurate about the fundamental, reducing coordination failure. We call this the ``wisdom of AI'' effect. At high $\rho$, the same correlation acts like a public signal: all AI agents respond identically to the common shock $\eta$, creating correlated withdrawal patterns that make bank runs self-fulfilling even at fundamentals where the bank is solvent. We call this the ``herding of AI'' effect. The transition occurs at $\rho_0$. Below it, AI stabilises. Above it, AI destabilises.

This non-monotonicity is absent from the \citet{danielsson2025} model. Their crisis threshold $\theta^*(\mu) = (c/b)[(1-\mu)p + \mu]$ is linear in $\mu$ because the AI speed advantage enters additively. Our model operates through information structure, not execution speed. The distinction matters: the two frameworks produce opposing risk assessments for the same market.

A numerical comparison illustrates the difference. With $\lambda = 0.5$, $\alpha_{\text{SC}} = 1$ (so $\rho_1^* = 0.5$), and $\rho_0 = 0.3$, Configuration A has many but diverse AI agents ($\mu = 0.8$, $\rho = 0.1$). The Danielsson-Uthemann model classifies it as dangerous; ours classifies it as safe ($\rho < \rho_1^*$). Configuration B has few but homogeneous AI agents ($\mu = 0.2$, $\rho = 0.7$). Their model classifies it as safe; ours classifies it as fragile ($\rho > \rho_1^*$). The two primitives produce opposing risk assessments.

\subsection{Welfare Analysis}

\begin{proposition}[Social Cost of AI Signal Correlation]\label{prop:welfare-channel1}
  In the Angeletos-Pavan (2007) framework applied to the bank-run coordination game, the social welfare loss from signal correlation is
  \begin{equation}\label{eq:welfare-loss}
    W_{\text{loss}}(\rho) = \frac{\alpha_{\text{SC}}^2}{(1-\alpha_{\text{SC}})^2} \cdot \lambda^2 \cdot \frac{\rho}{\tau},
  \end{equation}
  where $\tau = 1/\sigma^2$ is the total signal precision. The welfare loss is linearly increasing in $\rho$ for given $\alpha_{\text{SC}}$ and $\lambda$. The private value of adopting correlated AI signals is positive: $V_{\text{priv}}(\rho) = (1-\rho)/\tau > 0$ (reduced idiosyncratic risk). For $\alpha_{\text{SC}} > 0$ and $\lambda > 0$, the socially optimal correlation $\rho^{\text{SO}}$ is strictly below $\rho_1^*$.

  \prooflabel{Proof sketch.} The \citet{angeletos2007} (Proposition 4) welfare result states that in a coordination game with strategic complementarity $\alpha$, agents place weight $\alpha/(1-\alpha)$ on common information relative to the efficient benchmark, creating a welfare loss proportional to the square of this over-weighting. The variance of the common action component is $\lambda^2 \cdot \rho\sigma^2$, giving the welfare loss~\eqref{eq:welfare-loss}. The private benefit of correlation is the reduction in idiosyncratic decision error: as $\rho$ increases, the private noise $(1-\rho)\sigma^2$ falls. The social cost exceeds the private benefit when $[\alpha_{\text{SC}}^2/(1-\alpha_{\text{SC}})^2] \cdot \lambda^2 > 1$, which holds for $\alpha_{\text{SC}}$ and $\lambda$ above threshold values. $\square$
\end{proposition}

Proposition~\ref{prop:welfare-channel1} identifies the coordination externality of AI adoption. Each agent benefits from reducing her idiosyncratic signal noise. But the correlated errors impose a negative externality on all other agents by making coordinated bank runs more likely. The externality is proportional to the square of the strategic complementarity parameter, creating a wedge between private and social incentives that grows with $\rho$.

The \citet{szkup2015} uniqueness conditions interact with the correlated signal structure. Szkup and Trevino provide sufficient conditions for uniqueness in global games with endogenous information acquisition. Their key assumption, that signal precision is independent of other agents' information choices, is violated at high $\rho$. The informativeness of the aggregate signal depends on $\lambda$, which is itself an equilibrium object when Channel 2 is operative. At $\rho > \rho_1^*$, the Hellwig mechanism directly restores multiplicity, providing a more fundamental violation than endogeneity of the information structure alone.

Channel 1 establishes that AI signal correlation creates coordination risk through an information-theoretic mechanism. The common shock $\eta$ functions as an endogenous public signal that can coordinate agents on a bank-run equilibrium. The crisis threshold rises, and the uniqueness selection breaks down.

But coordination failure alone does not explain why the market's information aggregation deteriorates. Even in the uniqueness region ($\rho < \rho_1^*$), agents are acting on increasingly similar signals, and the price system reflects less and less genuine diversity of opinion. The question is: what happens to the information that prices are supposed to reveal? Channel 2 answers this question by showing that rising $\rho$ degrades revelatory price efficiency, collapsing the information ecosystem from within.
% --- Clarity pass complete: 2026-03-10 ---
% --- Flow pass complete: 2026-03-11 ---
% --- Technical audit complete: 2026-03-11 ---
